\begin{claim}{D.47}{Proved}{Conjugator coordinates at the binary refinement}
Use the specified quotient
$\widetilde Q=S^{19}\times C_2^2$, $S=\mathrm{SL}_2(\F_5)$,
of D.45, and its $324\,000\,000$ refined pairs $(t,f)$.
Every one of these pairs satisfies D.44's finite
first-symmetry coboundary condition
\[
       h^{-1}\alpha(h)=f,\qquad h\theta(h)=1,
                         \qquad h\in S^{19}.             \tag{D.47a}
\]
Exactly $10\,125\,000$ of the pairs satisfy the two
finite conjugator equations of D.35 at this quotient.
The count stays the same after requiring both actual
invertibility of the finite substitution and both
Harbater--Schneps symmetries for the conjugating pair.
There are $5\,062\,500$ pairs with admissible unit
parameters congruent to a square modulo five, and
$5\,062\,500$ with parameters congruent to a nonsquare.
These two sets are disjoint.

Thus $313\,875\,000$ refined coordinate pairs lie
outside the image of the standard full $GT_0$
involution conjugacy class.  For each of D.38's
$2\,531\,250$ conjugator-compatible base tuples,
exactly four of its $64$ refined $f$-coordinates
pass this test; for each of the other
$2\,531\,250$ base tuples, none does.
These are statements about the exact normalized
coordinates: D.45's $64$ coordinates for a fixed
$t$ define the same finite automorphism.
No surviving pair is asserted to have a full
involution or conjugator lift.
\end{claim}
\begin{proof}
Write $X_i,Y_i,Z_i,\widetilde t_i$ and $\alpha$ for
the matrices and actual automorphism fixed in D.45.
Since $S$ has exponent dividing $60$, the power
parameter in D.35 factors through
$(\Z/60)^\times$.  For the second Harbater--Schneps
symmetry we retain the finer, explicit parameter map
\[
 a\in(\Z/120)^\times\longmapsto
       \left(a\bmod60,\ m=(a-1)/2\bmod60\right).
                                                        \tag{D.47b}
\]
It is well-defined, and every such $a$ is the
reduction of a unit in $\widehat\Z$.
Changing the representative of $a$ by $120$ changes
$m$ by $60$, so all subsequent equations are
independent of the choice of integer representative.

The image of $F'$ is exactly $S^{19}$ by D.45.
For fixed $(t,f)$ the first conjugator equation
therefore separates into the nineteen exact equations
\[
 h_i^{-1}\widetilde t_i h_i\widetilde t_i^{-1}=f_i,
                                  \qquad h_i\in S.       \tag{D.47c}
\]
Each has exactly four solutions.  The program
verifies this by enumeration of all $120$ matrices.
It also follows from the order-four conjugacy class
in $S$: $\widetilde t_i$ has square $-I$, and the
norm equation in D.45 gives
$(f_i\widetilde t_i)^2=\widetilde t_i^2=-I$.
The thirty matrices with square $-I$ form one
conjugacy class with centralizer of order four,
as the same matrix enumeration verifies.

For each parameter $a\bmod60$ and each of these
four matrices, the second equation of D.35 is
tested by the exact matrix conjugacy condition
\[
      X_i^a h_iY_i^a h_i^{-1}
                       \sim_S (X_iY_i)^a.               \tag{D.47d}
\]
Indeed the $C_2^2$ coordinates of both sides are
$e_x+e_y$, since $a$ is odd, and that direct factor
is central.  Conjugacy in $\widetilde Q$ is therefore
exactly these nineteen independent matrix conjugacies.
The program computes every conjugacy class in $S$
by conjugation by all $120$ matrices; (D.47d) is
tested against those complete classes.

The finite invertibility test is explicit as well.
Reduce the ordered pair
$(X_i^a,h_iY_i^a h_i^{-1})$ to $A_5$ and test whether
it belongs to the $i$th recorded $\Aut(A_5)$ orbit
of generating pairs in D.38.
This is sufficient: its unique component
automorphism lifts to an automorphism $\gamma_i$
of $S$, and the two central differences from
$\gamma_i(X_i),\gamma_i(Y_i)$ are corrected by
characters of $C_2^2$, exactly as in (D.45d)--(D.45e).
The resulting automorphism of $\widetilde Q$ has
the prescribed two generator images.
It is necessary for a finite pure automorphism
with these images: its reduction preserves the
three specified power-conjugacy classes, and
D.38's distinct passports force it to preserve
each simple factor.  Its pair must therefore be
in that recorded orbit.  This tests the actual
finite substitution, with its characteristic
quotient map retained.

For completeness the program also imposes both
symmetries on the conjugating pair, not just its
three peripheral images.
Put $k=h^{-1}$ and use $m$ from (D.47b).
The equations tested are
\[
       h\theta(h)=1,\qquad
       \omega^2(k)Z^m\omega(k)Y^m kX^m=1.
                                                        \tag{D.47e}
\]
All $h,k,f$ coordinates lie in the perfect factor
$S^{19}$; the actions of $\theta,\omega$ there
are the specified lifted factor maps of D.45.
The $C_2^2$ coordinate of the second equation is
$m(e_x+e_y+e_z)=0$.
Thus these are exact equations at the indicated
finite quotient, including the half-exponent
required by relation (II).

The first equation of (D.47e) is exhausted without
assuming independent choices where factors are
exchanged.  If
$\theta(v)_i=G_i v_jG_i^{-1}$, a choice of $h_i$
forces
\[
                       h_j=G_i^{-1}h_i^{-1}G_i.         \tag{D.47f}
\]
For a fixed factor it must satisfy the same equation
with $j=i$.  The program enumerates one of the
four candidates from (D.47c) for each orbit of
$\theta$, uses (D.47f), and checks membership in
the corresponding candidate sets.
It verifies that (D.47a) is soluble for all the
refined pairs.  For the full conjugator test it
then checks (D.47d), finite invertibility, and the
second equation of (D.47e) for every remaining
choice and all $32$ units modulo $120$.

Here is the entire factorization of the result.
Let $U_+$ and $U_-$ denote the subsets of
$(\Z/120)^\times$ whose residues modulo five are
respectively squares and nonsquares; each has
$16$ elements.
For each of D.45's six joint factor orbits, the
table gives the number of refined pairs having
the specified set of admissible parameters.
\[
\begin{array}{c|r|r|r|r}
 \text{factor orbit}&U_+&U_-&U_+\cup U_-&\varnothing\\ \hline
 \{0,2,10\}&15&15&0&0\\
 \{1,9,17\}&15&15&0&0\\
 \{3,5,6,7,11,14\}&30&30&0&60\\
 \{4,12,18\}&15&15&0&0\\
 \{8,13,15\}&10&10&0&0\\
 \{16\}&0&0&5&0.
\end{array}                                               \tag{D.47g}
\]
The admissible global parameter set is the
intersection of the six local sets, because
the parameter is common while the factor
equations have the exact orbit decomposition
of D.45.
It is nonempty precisely when the first five
nonempty choices all select the same one of
$U_+,U_-$.
Each choice gives
$15\cdot15\cdot30\cdot15\cdot10\cdot5
 =5\,062\,500$ pairs.
Their sum is $10\,125\,000$, and subtraction
from D.45's $324\,000\,000$ gives the stated
$313\,875\,000$ failures.
The program records the same table when only
(D.47c)--(D.47d) are imposed, and when finite
invertibility and the first symmetry are added.
Thus the asserted equality of these counts is
itself checked at every factor orbit.
It also records the counts for each base tuple
before multiplying over the six orbits.  Their
Cartesian product gives the exact histogram
\[
 \begin{array}{c|r}
 \text{admissible }f\text{-coordinates per base tuple}
          &\text{number of base tuples}\\ \hline
 0&2\,531\,250\\
 4&2\,531\,250.
 \end{array}                                             \tag{D.47h}
\]

Every full conjugator in D.35 reduces to solutions
of (D.47c)--(D.47d), and D.35's equality of the
$P$ and $GT_0$ conjugator sets also supplies
(D.47e).
Consequently failure of the displayed finite
test excludes that exact coordinate pair from
the image of the standard full conjugacy class.
The executed implementation is
\path{code/part_a/pure_involution_binary_conjugacy.py}.
Its output retains all unit sets, the six local
tables, the per-base histogram, and the separate
first-symmetry counts.  The proof uses only this
finite refinement and makes no assertion that
one of its surviving pairs lifts through all
characteristic refinements of $F$.
\end{proof}
